% ============================================================================
%  CAPÍTULO II: OBJETIVOS, ALCANCE Y LÍMITES
% ============================================================================

\section{Objetivos}

\subsection{Objetivo general}

Migrar de forma íntegra, segura y verificable la base de datos de créditos y
ahorros de la Cooperativa de Ahorro y Crédito Horizonte PyME S.R.L. desde
MySQL 5.7 hacia PostgreSQL 16, sobre un entorno Linux reproducible con Docker
Compose, aplicando endurecimiento, roles de mínimo privilegio, enmascaramiento
de datos sensibles y respaldos cifrados.

\subsection{Objetivos específicos}

\begin{enumerate}
  \item Diagnosticar el sistema de origen: versión, tamaño, tablas, vistas,
        índices, usuarios, privilegios y riesgos de seguridad.
  \item Identificar y documentar las incompatibilidades de tipos y semántica
        entre MySQL 5.7 y PostgreSQL 16, con su regla de conversión.
  \item Diseñar el esquema destino con tipos nativos, restricciones
        \texttt{CHECK}, claves foráneas, índices y vistas equivalentes.
  \item Implementar un ETL determinista que preserve identificadores,
        integridad y zona horaria, y que verifique conteos origen--destino.
  \item Endurecer PostgreSQL 16 con TLS/SSL obligatorio,
        \texttt{scram-sha-256}, revocación de \texttt{PUBLIC}, roles granulares
        y registro de conexiones y DDL.
  \item Enmascarar documentos, teléfonos, correos y números de cuenta en una
        vista de auditoría, sin exponer las tablas base.
  \item Generar respaldos cifrados con AES-256-CBC, calcular su huella
        SHA-256 y probar la restauración completa.
  \item Automatizar la auditoría, la migración y la verificación en scripts
        reproducibles, con evidencias reales en \texttt{evidence/}.
\end{enumerate}

\section{Alcance y límites}

\subsection{Alcance}

El alcance incluye el entorno de laboratorio completo: los servicios Docker
de origen y destino, el esquema y los datos sintéticos del origen, el esquema
destino, el ETL, los scripts de auditoría, migración, respaldo y verificación,
las evidencias generadas y la documentación técnica del proceso de migración
y endurecimiento. Se
cubren las 8 tablas, las 2 vistas, los 19 índices, el usuario legacy, las
incompatibilidades de tipos, el hardening del motor destino, los roles, el
enmascaramiento, el cifrado en tránsito y en reposo, y las pruebas de
integridad y rendimiento básico.

\subsection{Límites}

Quedan fuera del alcance la puesta en producción real, la migración de datos
históricos de la cooperativa, la modificación del código de la aplicación, la
implementación de alta disponibilidad con réplicas o \emph{failover}, la
instalación de un \emph{pool} de conexiones como PgBouncer, la certificación
formal de cumplimiento normativo y la operación 24×7. Los datos utilizados
son ficticios y deterministas; no se manipuló información real de personas.
El entorno se ejecuta en un solo host Linux y no representa una arquitectura
de producción multizona. Estas exclusiones se documentan para que el lector
no confunda el laboratorio con un despliegue productivo.
